\section{Certification à profondeur \texorpdfstring{\(10^4\)}{10^4} des développements coinductifs de \texorpdfstring{\(\pi\), \(e\) et \(\varphi\)}{pi, e et phi}}
\label{sec:truncamiento-10000-pi-e-phi-rev9}

\noindent\textbf{Portée mathématique.}
Le théorème uniforme précédent quantifie sur tout entier positif \(N\) et
construit l'objet infini comme limite inverse de ses cylindres compatibles.
Ce chapitre exécute le cas \(N=10^4\) comme certificat fini particulièrement
exigeant ; \(10^4\) n'est pas une borne du générateur, de même qu'une exécution
à un million de chiffres serait une autre instance du même algorithme et non une
modification de sa loi.

La présente exécution fixe une profondeur de calcul et vérifie que le
cylindre obtenu est contenu dans une unique cellule décimale de dix mille
positions. Le plus petit nombre de blocs compatible avec la monodromie nonadique
et avec
\[
  3^{6B}>10^{10000}
\]
est
\[
  B=3501=389\cdot9,
  \qquad 6B=21006\ \text{trits}.
\]

Pour chacun des trois canaux, les \(729\)
sous-cylindres sont parcourus exhaustivement à chaque étape et l'on conserve l'unique sous-cylindre compatible avec l'encadrement
rationnel du lecteur. Le calcul produit
\[
 \begin{array}{c|r}
 \text{grandeur calculée par canal}&\text{valeur}\\
 \hline
 \text{partitions complètes }729\to1&3501\\
 \text{trits transportés}&21006\\
 \text{paires de même phase }(K,K+9)&3492\\
 \text{tours nonadiques complets}&388
 \end{array}
\]
et ne nécessite pas de nonade supplémentaire. Les \(396\) premiers blocs
coïncident exactement avec le témoin précédent à mille chiffres.

\Needspace{28\baselineskip}
Les lecteurs de clôture, de propagation et d'auto-échelle sont évalués au moyen
d'inégalités rationnelles exactes. L'accélération du calcul réorganise les
sommes par l'arithmétique entière et n'ouvre ni bibliothèques de constantes ni préfixes
décimaux cibles. Ce n'est qu'après avoir écrit les sorties que l'on exécute un
second calcul indépendant. On vérifie, dans cet ordre, la
génération rationnelle des trois canaux, la composition G9--T1 de
\(\alpha_{\rm HMT}\) et la coïncidence des dix mille chiffres avec les
développements de référence consultés seulement à la fin. Les trois comparaisons
sont exactes sur les dix mille positions calculées. Cette vérification finie
ne remplace pas l'argument uniforme de prolongation à profondeur arbitraire.

La démonstration étendue précise aussi le sens du retour. Dans les
\(3492\) paires par canal, la phase est conservée et l'état
complet change toujours, car la profondeur, le préfixe et le cylindre avancent. Une projection locale
réduite peut se répéter dans certaines paires sans que l'état TPK revienne. Cette
distinction remplace l'affirmation plus forte d'apériodicité locale par la
propriété effectivement démontrée : monodromie de phase avec mémoire de l'état
enrichi.

Les quatre développements complets sont imprimés dans
\cref{app:desarrollos-10000-rev9}. Les récurrences, les intervalles
rationnels et les comptages précédents déterminent le calcul ; les développements
imprimés sont leurs troncatures finies et permettent de comparer chaque chiffre sans
devenir des prémisses du générateur.
